% s02 The model, King's cube and the standing lemmas
\section{The model, King's cube and the standing lemmas}

\subsection{Couplings, lattices and the action}

Throughout, $e>0$ and $\lambda,m,\lmu>0$ are the charge, the quartic coupling, the bare scalar mass and the photon mass; no further condition on the couplings is imposed, since every smallness condition of \cite{K1,K2,Ba1,Ba2,Ba3,Ba4} is a cap on the running couplings at a stopping scale chosen after the point $(e,\lambda,m,\lmu)$ and after the datum (the section on the coupling range). The block factor is $5$: the spacings are $\eps_k\defeq5^{-k}$, the lattices $\Lk\defeq\eps_k\Z^3$ with the standard frame $(e_\alpha)$, and $M\in\N$ is King's large-block integer, fixed once and for all. King's cubes are $\CJ\defeq[-MJ,MJ)^3$, $J\in\N$, and $\XJ\defeq\CJ\cap\Lk$ (the spacing understood from the context). A \emph{King torus} is $\R^3/(\ell\Z)^3$ with $\ell\in M\N$ even; $\Lk$ descends to it, with $\ell\,5^k$ sites per side.

The scalar field is $\phi(x)=r_xe^{\mathrm i\vartheta_x}\in\C$ and the photon field a real $A_\alpha(x)$ for every site $x$ and direction $\alpha$; the bond $\langle x,x+\eps e_\alpha\rangle$ carries $A_\alpha(x)$ and the reversed bond $-A_\alpha(x)$. King's action \cite[(2.5)]{K2} at spacing $\eps$ is
\begin{equation*}
\begin{aligned}
S&=\half\sum_{\alpha,\beta}\sum_x\eps\,\big(A_\alpha(x+\eps e_\beta)-A_\alpha(x)\big)^2+\frac{\lmu^2}2\sum_x\eps^3|A(x)|^2
+\half\sum_\beta\sum_x\eps\,\big|e^{\mathrm ie\eps A_\beta(x)}\phi(x+\eps e_\beta)-\phi(x)\big|^2\\
&\quad+\sum_x\eps^3\Big[\half\big(m^2+\dm_\eps\big)r_x^2+\lambda r_x^4\Big],
\end{aligned}
\end{equation*}
the photon in Feynman form with mass $\lmu$, and the hopping term of a bond $\langle x,y\rangle$, $y=x+\eps e_\beta$, equal to $\frac\eps2(r_x^2+r_y^2)-\eps r_xr_y\cos(\vartheta_x-e\eps A_\beta(x)-\vartheta_y)$: a \emph{diagonal part} and a \emph{cosine part}. The counterterm $\dm_\eps$ is King's free-boundary counterterm \cite[p.~330]{K2}, fixed by his two-point graphs of order at most four in $(e,\lambda)$ at zero momentum with the propagators of the infinite lattice $\eps\Z^3$: it is \emph{one number per spacing}, the same at every site of every box and torus. The Wick constant is the free one, $\sW_\eps\defeq(-\Delta_{\eps\Z^3}+m^2)^{-1}(0,0)$. This \emph{free convention} is used on cubes and tori alike unless said otherwise; King's own torus convention appears only inside the torus inputs of the Wick leg.

\subsection{The states}

\begin{definition}[Neumann--Dirichlet boxes, killed sites, the comparison cube]\label{def:s02:acube}
For a finite set $X\subset\Lk$ the \emph{Neumann--Dirichlet state} $\langle\cdot\rangle^{\ND}_{k,X}$ is the normalised measure $Z^{-1}e^{-S}\prod_{x\in X}d^2\phi(x)\prod_{x\in X,\alpha}dA_\alpha(x)$ with $\phi=0$ at every site outside $X$ (Dirichlet condition: a bond leaving $X$ contributes only its diagonal part $\frac\eps2r_x^2$ at its end $x\in X$), photon gradient terms only for pairs $x,x+\eps e_\beta$ both in $X$, and the photon mass term on $X$ (Neumann condition). For $U\subset X$ the state \emph{with the matter killed on $U$} replaces the single-site measure $d^2\phi(x)$ times its weight by the point mass at $\phi(x)=0$ for $x\in U$; every hopping term with an endpoint in $U$ then reduces to its diagonal part at the other endpoint, and the photon variables are unchanged. The \emph{comparison cube} is $\NDa_{k,J}\defeq\langle\cdot\rangle^{\ND}_{k,\XJ}$, written $\ava{\cdot}$: the scalar is alive on all of $\XJ$, with Dirichlet planes $x_\mu=-MJ-\eps_k$ and $x_\mu=MJ$.
\end{definition}

\begin{definition}[King's RG cube]\label{def:s02:Kcube}
Let $\alive^K_J\defeq\{x\in\XJ:\ x_\mu<MJ-\eps_k\ \forall\mu\}$ and $\TJ\defeq\XJ\setminus\alive^K_J$, the three top face layers. \emph{King's RG cube} $\NDK_{k,J}$, written $\avK{\cdot}$, is the Neumann--Dirichlet state on $\XJ$ with the matter killed on $\TJ$: the photon variables $A_\alpha(x)$, $x\in\XJ$, their gradient terms for pairs in $\XJ$ and their mass term on $\XJ$ are those of the comparison cube; the scalar is alive on $\alive^K_J$ and $\phi=0$ on every other site of $\Lk$, in particular on $\TJ$. Its Dirichlet planes are $x_\mu=-MJ-\eps_k$ and $x_\mu=MJ-\eps_k$, at separation $2MJ$; the Neumann image planes of the photon, $x_\mu=-MJ-\eps_k/2$ and $x_\mu=MJ-\eps_k/2$, are the same in both conventions.
\end{definition}

The two conventions differ exactly by the matter on $\TJ$; the photon Gaussian is the same. King's RG cube is the box \cite[(3.3)--(3.4)]{K2} with the Dirichlet planes of \cite[(3.9)]{K2} and the photon reflection planes of \cite[(3.7)]{K2}, the geometry of his renormalisation group; his Theorem 3.2 and his modified box \cite[(3.38)]{K2} are written in the comparison convention. The two families are interlaced in Section~4; the existence theorems of this paper are posed on $\NDK$, while the rotation part and the cell lemmas of the Wick leg are written for the comparison family. The \emph{periodic state} $\langle\cdot\rangle_{k,T}$ on a King torus $T$ has the action summed over the finite graph $\Lk/(\ell\Z)^3$, in the free convention.

\paragraph{BFS measures.} A \emph{BFS measure} on a finite set $V$ of sites is a probability measure
\[
Z^{-1}\exp\Big(\sum_{\{x,y\}\in B}c_{xy}\,r_xr_y\cos\big(\vartheta_x-\vartheta_y-\ell_{xy}(A)\big)\Big)\prod_{x\in V}d\rho_x(r_x)\,d\vartheta_x\;d\gamma(A),
\]
with $B$ a set of pairs of sites, $c_{xy}\ge0$, each $\rho_x$ a finite measure on $[0,\infty)$ with $\int e^{\alpha r^2}d\rho_x<\infty$ for all $\alpha\ge0$ (point masses at $0$, that is killed sites, included), $\gamma$ a non-degenerate centred Gaussian measure on $\R^m$ (the photon variables), and linear functionals $\ell_{xy}$. Every state above is a BFS measure: $c_{xy}=\eps$ for the bonds present, the single-site measures contain $e^{-\eps^3[\frac12(m^2+\dm_\eps)r^2+\lambda r^4]}$ and the diagonal parts $\frac\eps2r_x^2$ of all bonds at $x$ (those leading out of the box or into a killed site included); the photon form is positive definite because $\lmu>0$, and the exponential moments are finite because $\lambda>0$. So is every measure obtained from such a state by multiplying some $c_{xy}$ or some photon gradient terms by $t\in[0,1]$, adding $s|\ell(A)|^2$ ($s\ge0$) to the action, multiplying single-site measures by $e^{\eps^3q(x)r^2}$ with real $q$, or killing sites. A BFS measure is \emph{of King type} if $c_{xy}\in[0,\eps]$ and each $\rho_x$ is a point mass at $0$ or has the density $r\,e^{-\eps^3[\frac12(m^2+\dm_\eps)r^2+\lambda r^4]-\eps^3w(x)r^2-\frac\eps2n_xr^2}$ with $w(x)\ge0$, $n_x\in\{0,\dots,6\}$; every state above, and every measure obtained from one by the operations just listed with $q\le0$, is of King type. Every state above is invariant under the \emph{charge conjugation} $(A,\vartheta)\mapsto(-A,-\vartheta)$, which changes the sign of every field strength and fixes $r$.

\subsection{Sources, schemes and the objects}

\begin{definition}[sources, schemes, classes, transport]\label{def:s02:sources}
A \emph{curvature source} is a real $2$-form $f\in C^\infty_c(\R^3;\Lambda^2)$, a \emph{scalar source} a real $h\in C^\infty_c(\R^3)$. For the plaquette $p=(x;\alpha<\beta)$ with centre $c(p)$ let $F(p)\defeq\eps^{-1}(A_\alpha(x)+A_\beta(x+\eps e_\alpha)-A_\alpha(x+\eps e_\beta)-A_\beta(x))$ and
\[
F_{\Lk}(f)\defeq\sum_p\eps^3f_{c(p)}(e_\alpha,e_\beta)F(p),\qquad
\sS_{\Lk}(h)\defeq\sum_x\eps^3h(x),\qquad
\Wick{h}\defeq\sum_x\eps^3h(x)\big(|\phi(x)|^2-2\sW_\eps\big),
\]
the sums over the plaquettes (corners in the box), respectively the sites, of the region. A \emph{scheme} is a real sequence $a=(a_k)$, $a_k$ depending only on $\eps_k$ and the couplings; $W^a(h)\defeq\sum_x\eps^3h(x)(|\phi(x)|^2-2\sW_\eps-a_k)$. For a finite family $\zeta=(z_j,f_j;w_j,h_j)_j$ with complex coefficients put $\Theta(\zeta)\defeq\exp(\sum_jz_jF(f_j)+\sum_jw_j\Wick{h_j})$ and, for a Neumann--Dirichlet state on a box $X$ (killed sites allowed),
\[
Z^a_{k,X}(\zeta)\defeq\Big\langle\exp\Big(\sum_jz_jF(f_j)+\sum_jw_jW^a(h_j)\Big)\Big\rangle^{\ND}_{k,X}
=e^{-\sum_jw_ja_k\sS_{\Lk}(h_j)}\big\langle\Theta(\zeta)\big\rangle^{\ND}_{k,X};
\]
$Z^a_{k,J}\defeq Z^a_{k,\XJ}$ for the comparison cube, $Z^{a,K}_{k,J}$ for King's RG cube, and $N^a_{k,X}(0,q)\defeq e^{-(a_k+2\sW_\eps)\sS(q)}\langle e^{\sum_x\eps^3q(x)r_x^2}\rangle^{\ND}_{k,X}$ for one real scalar source. At $h=0$ the multiplier is $1$ and $a$ is dropped: $Z_{k,J}(z)\defeq\ava{e^{\sum_jz_jF(f_j)}}$, $Z^K_{k,J}(z)\defeq\avK{e^{\sum_jz_jF(f_j)}}$ for $z\in\C^n$. The \emph{cone} is $\cK\defeq\{(f,g):f\text{ a real $2$-form},\ g\ge0\}$ with the observable $\Theta(\mathrm if,g)$ (\emph{class $\cI$}); \emph{class $\cD$} consists of the observables $\Theta(f,-g)$, $(f,g)\in\cK$. Transport: $(O\cdot h)(x)\defeq h(O^{-1}x)$, $(O\cdot f)_x(u,v)\defeq f_{O^{-1}x}(O^{-1}u,O^{-1}v)$ for $O\in\Oth$; $\tau_bh\defeq h(\cdot-b)$, $(\tau_bf)_x\defeq f_{x-b}$. For $\nu\in\{1,2,3\}$, $c\in\R$ let $\rho=\rho_{\nu,c}$, $x\mapsto x-2(x_\nu-c)e_\nu$; $\theta h\defeq h\circ\rho$, $\theta f\defeq\rho^*f$, $(\rho^*f)_x(u,v)\defeq f_{\rho x}(\rho u,\rho v)$ (the components $f_{\mu\nu}$, $\mu\ne\nu$, change sign); $\theta\zeta^*\defeq(\bar z_j,\rho^*f_j;\bar w_j,h_j\circ\rho)_j$, and $\zeta\sqcup\zeta'$ is the union of two families. The \emph{amplitude} is $\amp(\zeta)\defeq\sum_j|\Re w_j|\,\norm{h_j}_\infty$. The \emph{reference scheme} is $z=a^\sharp$, $a^\sharp_k\defeq e^2X_{\eps_k}$, $X_\eps$ the coefficient of $e^2$ in the perturbative expansion of $\langle|\phi(0)|^2\rangle$ on $\eps\Z^3$ at $\lambda=0$ with the order-$e^2$ part of $\dm_\eps$ inserted.
\end{definition}

By charge conjugation $\langle\Theta(\mathrm if,g)\rangle=\langle\cos F(f)\,e^{\Wick{g}}\rangle$ is real for every state above, and $Z^a_{k,X}(-z,w)=Z^a_{k,X}(z,w)$.

\begin{remark}[schemes]\label{rem:s02:scheme}
The Wick leg is stated for the reference scheme $z$, or for an admissible scheme $a$ (\cref{def:s03:adm}: $z_k-a_k$ converges, hence is bounded). If the fixed-cube limits exist in $(0,\infty)$ for both $a$ and $z$ at a real scalar source, the scheme is admissible: for $0\le h\in C^\infty_c$ with $\int h>0$ supported well inside some $\CJ$, $Z^z_{k,J}(0,h)/Z^a_{k,J}(0,h)=e^{-(z_k-a_k)\sS_{\Lk}(h)}$ converges to a number in $(0,\infty)$, and $\sS_{\Lk}(h)\to\int h>0$. What the Wick leg uses is only $\sup_k|z_k-a_k|<\infty$.
\end{remark}

\begin{definition}[norms and the tube]\label{def:s02:tube}
For a family with sources $f_1,\dots,f_n$, $h_1,\dots,h_p$ supported in $Q_R\defeq[-R,R]^3$ and $\zeta=(z,w)\in\C^{n+p}$ put $g(\zeta)\defeq\sum_i(\Re z_i)f_i$ and $k(\zeta)\defeq\sum_j(\Re w_j)h_j=k_+-k_-$ (pointwise positive and negative parts). Lattice norms: $\norm{g}_{2,k}^2\defeq\sum_p\eps_k^3g_{c(p)}(e_\alpha,e_\beta)^2$ over the plaquettes of $\Lk$, $\norm{q}_{1,k}\defeq\sum_{x\in\Lk}\eps_k^3|q(x)|$. Continuum norms: $\norm{g}_2^2\defeq\sum_{\alpha<\beta}\norm{g_{\alpha\beta}}_{L^2}^2$; $\wnorm{q}\defeq\sup_x(1+|x|^2)^2|q(x)|$; $\fournorm{q}\defeq\sup_{x,|\gamma|\le4}(1+|x|^2)^2|D^\gamma q(x)|$ and $\fournorm{g}\defeq(\sum_{\alpha<\beta}\fournorm{g_{\alpha\beta}}^2)^{1/2}$, as in \cite{SymPhase}. For $\sigma,H>0$ the \emph{tube} is
\[
\cD_{\sigma,H}\defeq\big\{\zeta\in\C^{n+p}:\ \norm{k_+(\zeta)}_\infty<\sigma/2,\ \ \norm{k_-(\zeta)}_\infty<H/2\big\}.
\]
Since $s\mapsto\norm{(\sum_js_jh_j)_\pm}_\infty=\sup_x(\pm\sum_js_jh_j(x))_+$ is a supremum of convex functions of $s=\Re w$, hence convex and continuous, $\cD_{\sigma,H}$ is the preimage under $\zeta\mapsto\Re w$ of the intersection of two open convex sublevel sets: open, convex, hence connected; it contains $0$ and every $\zeta$ with purely imaginary $w$; $\norm{k(\zeta)}_\infty<\min(\sigma,H)/2$ implies $\zeta\in\cD_{\sigma,H}$; its real points form an open subset of $\R^{n+p}$. The tube of the Wick leg is $\Dtube$, with $\sigma_+$ and $H$ fixed there.
\end{definition}

\paragraph{The objects.} For real $2$-forms $f_1,\dots,f_n$ let $J_0(f)$ be the least $J$ such that every support lies at $\ell^\infty$-distance $>1$ from $\R^3\setminus\CJ$. King's \emph{cube limits} and \emph{printed state} are
\[
\ZKJ(z)\defeq\lim_{k\to\infty}Z^K_{k,J}(z)\quad(J\ge J_0(f)),\qquad
\ZKinf\Big(\sum_jz_jf_j,0\Big)\defeq\lim_{J\to\infty}\ZKJ(z),
\]
whenever the limits exist: their existence at real $z$ is \cref{hyp:s03:EC}, and their existence at every $z\in\C^n$, entire in $z$, is derived from it in the curvature-sector section. The object is scheme-free. For general sources in a scheme $a$ the \emph{tube object} is $G_J(\zeta)\defeq\lim_kZ^a_{k,J}(\zeta)$, $\Ginf(\zeta)\defeq\lim_JG_J(\zeta)$, $\zeta\in\Dtube$, whenever these limits exist (the Wick-leg sections); where the full object $Z^a_\infty\defeq\lim_J\lim_kZ^a_{k,J}$ of the optional extension (\cref{thm:s23:FormA}) exists it agrees with $\Ginf$ on the tube, being the same iterated limit.

\subsection{Elementary facts}

\begin{lemma}[elementary facts]\label{lem:s02:E}
\emph{(E1)} For $\eps\le1$ and any translate $\Gamma$ of $\eps\Z^3$: $\sum_{x\in\Gamma}\eps^3(1+|x|^2)^{-4}\le C_0\defeq(5/2)^4\pi^2/8<49$ and $\sum_{x\in\Gamma}\eps^3(1+|x|^2)^{-2}\le C_0'\defeq(5/2)^2\pi^2<62$.
\emph{(E2)} $\norm{q}_{L^1}\le\pi^2\wnorm{q}$, $\norm{q}_{L^2}\le\frac{\pi}{2\sqrt2}\fournorm{q}$, $\norm{q}_\infty\le\wnorm{q}\le\fournorm{q}$, $\norm{g}_2\le\frac{\pi}{2\sqrt2}\fournorm{g}$; $\wnorm{\cdot}$ is a norm, so $\wnorm{k(\zeta)}\le\sum_j|\Re w_j|\wnorm{h_j}$ and $\wnorm{k_\pm}\le\wnorm{k}$.
\emph{(E3)} Sampling is evaluation at points: it commutes with pointwise positive and negative parts, and a sampled sup is at most the continuum sup, so $\zeta\in\cD_{\sigma,H}$ for the continuum sources implies the same bounds for the sampled ones. For continuous compactly supported $g,q$: $\norm{g}_{2,k}\to\norm{g}_2$, $\norm{q}_{1,k}\to\norm{q}_{L^1}$, and for every $k$: $\norm{g}_{2,k}^2\le C_0\fournorm{g}^2$, $\norm{q}_{1,k}\le C_0'\wnorm{q}$.
\emph{(E4)} $n^{n/2}\le e^{n/2}(n!)^{1/2}$.
\end{lemma}

\begin{proof}
(E1) For $y$ in the cube of side $\eps$ centred at $x\in\Gamma$, $|x-y|^2\le\frac34$ and $|y|^2\le2|x|^2+2|x-y|^2\le2|x|^2+\frac32$, so $1+|y|^2\le\frac52(1+|x|^2)$ and $(1+|x|^2)^{-r}\le(5/2)^r(1+|y|^2)^{-r}$. Integrate over the cube and sum over the disjoint cubes: $\sum_{x\in\Gamma}\eps^3(1+|x|^2)^{-r}\le(5/2)^r\int_{\R^3}(1+|y|^2)^{-r}dy$, with $\int(1+|y|^2)^{-4}=4\pi\cdot\frac\pi{32}=\frac{\pi^2}8$ and $\int(1+|y|^2)^{-2}=4\pi\cdot\frac\pi4=\pi^2$. (E2) $\norm{q}_{L^1}\le\wnorm{q}\int(1+|y|^2)^{-2}$, $\norm{q}_{L^2}^2\le\fournorm{q}^2\int(1+|y|^2)^{-4}$; the rest is immediate. (E3) Riemann sums for the convergences; the plaquette centres of one orientation form a translate of $\eps_k\Z^3$, so (E1) bounds $\sum_{p:(\alpha,\beta)}\eps_k^3g_{c(p)}(e_\alpha,e_\beta)^2$ by $C_0\fournorm{g_{\alpha\beta}}^2$, and the three orientations sum to $C_0\fournorm{g}^2$; the last bound is (E1) with $r=2$. (E4) is $n^n\le e^nn!$.
\end{proof}

\paragraph{Lattice-to-continuum.} On a box $X$ let $\ip AB_k\defeq\sum_{x\in X}\sum_\alpha\eps^3A_\alpha(x)B_\alpha(x)$, and for a real $2$-form $g$ whose non-vanishing plaquettes have their corners in $X$ define $\dstar_kg$ by $F_{\Lk}(g)=\ip{\dstar_kg}{A}_k$. Each component $(\dstar_kg)_\gamma$ is a backward difference quotient of the sampled components $g_{\gamma\beta}$, $\beta\ne\gamma$; for smooth compactly supported $g$ it converges uniformly to $(\dstar g)_\gamma$, up to the orientation sign convention of the components, at the rate $O(\eps_k)$ of Taylor's theorem. Hence
\begin{equation}\label{eq:s02:riemann}
\norm{\dstar_kg}_{2,k}\to\norm{\dstar g}_{L^2},\qquad \sup_k\norm{\dstar_kg}_{2,k}<\infty,\qquad \norm{g}_{2,k}\to\norm{g}_2\qquad(k\to\infty).
\end{equation}

\subsection{Correlation inequalities}

For a BFS measure, a real linear functional $F$ of $A$ and $g\ge0$ on $V$ put $\Theta_{F,g}\defeq e^{\mathrm iF}e^{\sum_xg_xr_x^2}$, $\Theta^-_{F,g}\defeq e^{F}e^{-\sum_xg_xr_x^2}$, and $\langle X;Y\rangle\defeq\langle XY\rangle-\langle X\rangle\langle Y\rangle$. Since $F(f)$ is a real linear functional of $A$ and $e^{\Wick{g}}$ a positive constant times $e^{\sum_x\eps^3g(x)r_x^2}$, class $\cI$ consists of multiples of observables $\Theta_{F,g}$ and class $\cD$ of multiples of observables $\Theta^-_{F,g}$.

\begin{lemma}[correlation inequalities]\label{lem:s02:BFS}
Let $\langle\cdot\rangle$ be a BFS measure, $F$ a real linear functional of $A$, $g\ge0$, $\Theta=\Theta_{F,g}$, $\Theta^-=\Theta^-_{F,g}$. Then $\langle\Theta\rangle$ is real, and
\begin{enumerate}[label=\textup{(\alph*)},nosep]
\item $\langle\Theta;|\ell(A)|^2\rangle\le0$ for every complex linear functional $\ell$ of $A$; hence adding $s|\ell(A)|^2$ to the action, $s\ge0$, does not decrease $\langle\Theta\rangle$, and deleting a photon gradient term does not increase it;
\item $\langle\Theta;r_xr_y\cos(\vartheta_x-\vartheta_y-\ell_{xy}(A))\rangle\ge0$; hence decreasing $c_{xy}$ does not increase $\langle\Theta\rangle$;
\item $\langle\Theta;r_x^2\rangle\ge0$; hence multiplying $\rho_x$ by $e^{-sr^2}$, $s\ge0$, does not increase $\langle\Theta\rangle$, and killing a set $U$ of sites does not increase it, provided $g=0$ on $U$ or the measure is of King type;
\item for $\Theta^-$ the three inequalities and the three monotonicity statements hold with the opposite signs;
\item if the pairs in $B$ and the Gaussian $\gamma$ do not couple the fields at the sites of $X\subset V$ (and the photon variables attached to them) to the others, the measure is a product, and for an observable depending only on the fields in $X$ the expectation equals the expectation in the $X$-factor.
\end{enumerate}
\end{lemma}

\begin{proof}
(a)--(c) are \cite[Thms.~6.1, 6.2(1),(2)]{BFS1}, whose hypotheses are those of a BFS measure: an arbitrary finite lattice, $c_{xy}\ge0$, finite single-site measures on $[0,\infty)$ with $\int e^{\alpha r^2}d\rho_x<\infty$ for all $\alpha\ge0$ (a point mass at $0$ is one), an arbitrary centred Gaussian with covariance $C\ge0$, and $g\ge0$ (so $g=0$ is included). Theorem 6.1 states $\langle P(r)\cos(n\cdot\vartheta+f\cdot A);Q(r)\cos(m\cdot\vartheta+h\cdot A)\rangle\ge0$ for polynomials $P,Q$ in the $r_x$ with positive coefficients; (a), (b) follow as in the printed proof of 6.2(1),(2): $e^{\sum g_xr_x^2}$ is a limit of such polynomials, $e^{\mathrm iF}$ may be replaced by $\cos F$ because the expectation and both observables are invariant under $(A,\vartheta)\mapsto(-A,-\vartheta)$ (which also gives the realness), and $|\ell(A)|^2=\ell_1(A)^2+\ell_2(A)^2$ with $\ell_r(A)^2=\lim_{t\downarrow0}t^{-2}[1-\cos(t\ell_r(A))]$; (c) is Theorem 6.1 with $Q=r_x^2$ and $m=0$, $h=0$. Each monotonicity statement follows by differentiating in the parameter, every interpolating measure being a BFS measure. Killing: if $g=0$ on $U$, set $c_{xy}=0$ on every bond at $U$ (not increasing, by (b)); the measure is then a product (e) of the matter at $U$, which carries no observable and cancels, and of the state with the matter killed on $U$ (the diagonal parts at the other ends of the bonds into $U$ are already in their single-site measures). For King type in general, multiply $\rho_x$, $x\in U$, by $e^{-s\eps^3r^2}$ (not increasing) and let $s\to\infty$ by Lemma~\ref{lem:s02:Kprime}.

(d) is stated in \cite[Thm.~6.2(3),(4)]{BFS1} for the photon and the hopping inequality with a sketched proof, and its $r^2$-inequality is \cite[(3.28)]{K2}, stated there without proof; we prove all three. Let $\mu$ be the unnormalised measure and $\omega'=(r',\vartheta',A')$ an independent copy of $\omega=(r,\vartheta,A)$; then $(\int d\mu)^2\langle\Theta^-;Y\rangle=\iint\Theta^-(\omega)(Y(\omega)-Y(\omega'))\,d\mu(\omega)d\mu(\omega')$. Write $r=u-v$, $r'=u+v$ (so $u_x\ge|v_x|$), $\vartheta=\vartheta_++\vartheta_-$, $\vartheta'=\vartheta_+-\vartheta_-$, $A=a_++a_-$, $A'=a_+-a_-$. Then $a_\pm=(A\pm A')/2$ are independent with the common centred Gaussian law of covariance $C/2$, and the radial law $\prod_xd\rho_x(u_x-v_x)d\rho_x(u_x+v_x)$ on $\{u\ge|v|\}$ is invariant under $v_x\mapsto-v_x$ at each site. With $\chi_{xy}\defeq\vartheta_{+,x}-\vartheta_{+,y}-\ell_{xy}(a_+)$ and $\psi_{xy}\defeq\vartheta_{-,x}-\vartheta_{-,y}-\ell_{xy}(a_-)$ the unprimed phase is $\chi+\psi$, the primed one $\chi-\psi$, and
\begin{align*}
r_xr_y\cos(\chi+\psi)+r'_xr'_y\cos(\chi-\psi)&=2(u_xu_y+v_xv_y)\cos\chi\cos\psi+2(u_xv_y+v_xu_y)\sin\chi\sin\psi,\\
r_xr_y\cos(\chi+\psi)-r'_xr'_y\cos(\chi-\psi)&=-2(u_xv_y+v_xu_y)\cos\chi\cos\psi-2(u_xu_y+v_xv_y)\sin\chi\sin\psi,
\end{align*}
$r_x^2-r_x'^2=-4u_xv_x$, $\ell(A)^2-\ell(A')^2=4\ell(a_+)\ell(a_-)$ for real $\ell$, and $\Theta^-(\omega)=e^{F(a_+)}e^{F(a_-)}e^{-\sum_xg_x(u_x^2+v_x^2)}e^{2\sum_xg_xu_xv_x}$. Expand $\exp(\sum_bc_b[\text{first line}])$ and $e^{2\sum g_xu_xv_x}$ in power series (absolutely convergent at fixed spacing). Every term of the integrand is a sign $\sigma$ times a monomial in $(u,v)$ with non-negative coefficient, times the even radial factor $e^{-\sum g_x(u_x^2+v_x^2)}$, times $V(\vartheta_+,a_+)V(\vartheta_-,a_-)$ with one real function $V$ (a product of $e^{F(\cdot)}$, of factors $\cos\chi_b$ or $\sin\chi_b$, and of $\ell(\cdot)$ for the photon slot). The photon integral of $V(\vartheta_+,a_+)V(\vartheta_-,a_-)$ factorises because $a_+,a_-$ are independent and identically distributed. The angular integral factorises by the following argument: per site $d\vartheta\,d\vartheta'=2\,d\vartheta_+d\vartheta_-$, and $[0,2\pi)^2$ in the $(\vartheta,\vartheta')$-variables is a fundamental domain of the lattice $L$ generated by $(\pi,\pi)$ and $(\pi,-\pi)$ in the $(\vartheta_+,\vartheta_-)$-plane, which contains $(2\pi\Z)^2$ with index $2$; $V(\vartheta_+,\cdot)V(\vartheta_-,\cdot)$ is $L$-invariant, because the simultaneous shift of $\vartheta_{+,x}$ and $\vartheta_{-,x}$ by $\pi$ at one site changes every $\chi_b$ and $\psi_b$ with $b\ni x$ by $\pi$ and leaves $\cos\chi_b\cos\psi_b$ and $\sin\chi_b\sin\psi_b$ unchanged; hence $\int_{[0,2\pi)^2}VV\,d\vartheta\,d\vartheta'=2\int_{L\text{-domain}}VV\,d\vartheta_+d\vartheta_-=\int_{[0,2\pi)^2}VV\,d\vartheta_+d\vartheta_-=(\int_0^{2\pi}V)^2$. The radial integral of a monomial vanishes if its degree in some $v_x$ is odd, by the $v\mapsto-v$ symmetry, and is non-negative otherwise, since $u_x\ge|v_x|$. So every term has the sign $\sigma$: $\sigma=-1$ for $Y=r_x^2$ (from $-4u_xv_x$) and for the hopping term (second line), $\sigma=+1$ for $Y=\ell(A)^2$ (for complex $\ell$, real and imaginary parts). Hence $\langle\Theta^-;r_x^2\rangle\le0$, $\langle\Theta^-;r_xr_y\cos(\cdot)\rangle\le0$, $\langle\Theta^-;|\ell(A)|^2\rangle\ge0$. Monotonicity and killing follow as in (a)--(c), with (d)'s hopping inequality in place of (b) and Lemma~\ref{lem:s02:Kprime}, which covers $\Theta^-$, for King type. (e) is clear.
\end{proof}

\begin{lemma}[current representation and Gaussian domination]\label{lem:s02:GD}
Let $\langle\cdot\rangle$ be a BFS measure whose Gaussian has covariance $C$. Integrating out the matter at fixed $A$, the matter integral is $\Psi(A)=\sum_qw(q)\cos\ip{\kappa(q)}A$, $q$ running over the integer-valued divergence-free currents on $B$, $w(q)\ge0$, $\sum_qw(q)<\infty$, $\ip{\kappa(q)}A\defeq\sum_{xy}q_{xy}\ell_{xy}(A)$. Let $\E^\sharp$ be the probability measure on currents with weights proportional to $w(q)e^{-\ip{\kappa(q)}{C\kappa(q)}/2}$. Then for $v\in\R^m$, $s\in\R$,
\[
\langle e^{s\ip vA}\rangle=e^{s^2\ip v{Cv}/2}\,\E^\sharp\cos\big(s\ip{Cv}{\kappa}\big)\le e^{s^2\ip v{Cv}/2}.
\]
\end{lemma}

\begin{proof}
Each bond factor expands as $e^{c_{xy}r_xr_y\cos t}=\sum_{q\in\Z}I_q(c_{xy}r_xr_y)e^{\mathrm iqt}$ with $I_q\ge0$, $I_{-q}=I_q$, $\sum_qI_q(c)=e^c$; after the radial integrals the multiple series converges absolutely, dominated by $\prod_xe^{\frac12\sum_yc_{xy}r_x^2}$, integrable against the single-site measures. Integrating the angles forces $\sum_yq_{xy}=0$ at every site (at a killed site $r_x=0$ and $I_q(0)=\delta_{q0}$, so the currents vanish on its bonds); pairing $q$ with $-q$ gives $\Psi$ with $w(q)\ge0$. For a centred Gaussian $\gamma$ with covariance $C$ and $a\in\C^m$, $\int e^{\ip aA}d\gamma=e^{\ip a{Ca}/2}$; with $a=sv\pm\mathrm i\kappa$ and the average of the two signs, $\int e^{s\ip vA}\cos\ip\kappa A\,d\gamma=e^{s^2\ip v{Cv}/2-\ip\kappa{C\kappa}/2}\cos(s\ip v{C\kappa})$. Sum over $q$ with the weights $w(q)$, divide by the same sum at $s=0$ (positive), and use $|\cos|\le1$.
\end{proof}

The lemma applies to every state of this paper at every $k$, $J$, $\lmu>0$, with or without a real scalar source absorbed in the single-site measures (the photon covariance is unchanged).

\begin{lemma}[killing]\label{lem:s02:Kprime}
Let $\langle\cdot\rangle$ be a BFS measure of King type, $U$ a set of alive sites, $\chi:U\to(0,\infty)$, and $\langle\cdot\rangle_{U,s}$ the measure with $\rho_x$ multiplied by $e^{-s\eps^3\chi(x)r_x^2}$ for $x\in U$. Let $\Theta$ be continuous with $|\Theta(A,\phi)|\le e^{\ell(A)}e^{C\sum_xr_x^2}$ for a real linear functional $\ell$ and $C\ge0$. Then $\langle\Theta\rangle_{U,s}\to\langle\Theta\rangle_U$ as $s\to\infty$, $\langle\cdot\rangle_U$ the measure with $\rho_x$ replaced by the point mass at $0$ for $x\in U$ (the state with the matter killed on $U$).
\end{lemma}

\begin{proof}
The unnormalised density is $e^{-Q(A)/2}dA\prod_{x\text{ alive}}w_x(r_x)r_xdr_xd\vartheta_x\prod_{b}e^{c_br_xr_y\cos(\cdot)}\prod_{x\in U}e^{-s\eps^3\chi(x)r_x^2}$, $Q$ positive definite, killed sites carrying no variables. Substitute $r_x=(s\eps^3\chi(x))^{-1/2}u_x$ for $x\in U$; the Jacobians cancel between numerator and denominator. \emph{Domination for $s\ge s_0$:} at $x\in U$, $w_x(r_x)\le e^{\eps^3|\dm_\eps|r_x^2/2}$ (the other exponents in $w_x$ are $\le0$), $|\Theta|$ contributes $e^{Cr_x^2}$, and each of the at most six bond weights at $x$ is at most $e^{\frac\eps2r_x^2+\frac\eps2r_y^2}$ ($0\le c_b\le\eps$); in the variable $u_x$ these factors together are $\le e^{u_x^2/2}$ once $s\ge s_0(\eps,m,\dm_\eps,C,\min\chi)$, so the $U$-factor is at most $e^{-u_x^2/2}u_x$, integrable on $[0,\infty)\times[0,2\pi)$; at an alive $y\notin U$ the factor is at most $w_y(r_y)e^{(C+3\eps)r_y^2}r_y$, integrable since $\lambda>0$; the photon factor $e^{\ell(A)}e^{-Q(A)/2}$ is integrable since $Q$ is positive definite. \emph{Pointwise limit:} at fixed $(u,\vartheta,A,r_{V\setminus U})$, $r_x\to0$ for $x\in U$, so $\Theta\to\Theta|_{\phi_U=0}$ by continuity, every bond weight with an endpoint in $U$ and every $w_x(r_x)$, $x\in U$, tends to $1$, and the $U$-factor tends to $e^{-u_x^2}u_x\,du_xd\vartheta_x$, whose integral cancels between numerator and denominator. By dominated convergence the ratio tends to $\int\Theta|_{\phi_U=0}d\mu_U/\int d\mu_U$, where $\mu_U$ has the cosines at $U$ removed, the weights $w_y$ (which already contain all diagonals at $y$) and the photon unchanged, and the point mass at $U$: that is $\langle\cdot\rangle_U$, with a positive denominator.
\end{proof}

\begin{lemma}[cutting]\label{lem:s02:cut}
Let $\langle\cdot\rangle_Y$ be a Neumann--Dirichlet box state on $\Lk$, possibly with killed sites, or a periodic state on a King torus, and let $X$ be a set of sites of $Y$ such that $X$ and its nearest neighbours lie in one fundamental cell (a clause in torus language: for a box state it has no content beyond $X\subset Y$). Let $\langle\cdot\rangle^{\ND}_X$ be the Neumann--Dirichlet state on $X$ whose single-site measures (killed sites and the diagonal parts of the bonds leaving $X$ included) are those of $\langle\cdot\rangle_Y$. Let $(f,g)\in\cK$ be such that every plaquette and every site at which $f$, respectively $g$, does not vanish lies in $X$. Then, at every lattice spacing,
\[
\langle\Theta(\mathrm if,g)\rangle^{\ND}_X\le\langle\Theta(\mathrm if,g)\rangle_Y,\qquad
\langle\Theta(f,-g)\rangle^{\ND}_X\ge\langle\Theta(f,-g)\rangle_Y,
\]
and the same for $Z^a$. The first inequality uses Lemma~\ref{lem:s02:BFS}(a),(b),(e) only; the second uses (d),(e). If no site of $X$ is killed, $\langle\cdot\rangle^{\ND}_X$ is the state of Definition~\ref{def:s02:acube}.
\end{lemma}

\begin{proof}
Delete from $Y$ every photon gradient term coupling a site of $X$ to a site outside $X$, and the cosine parts of all hopping terms between $X$ and its complement, keeping their diagonal parts $\frac\eps2(r_x^2+r_y^2)$. By Lemma~\ref{lem:s02:BFS}(a),(b) this does not increase the first expectation, by (d) it does not decrease the second, and every intermediate measure is a BFS measure. The result is a product (Lemma~\ref{lem:s02:BFS}(e)) whose $X$-factor is $\langle\cdot\rangle^{\ND}_X$: the diagonal parts of the cut bonds are exactly the Dirichlet terms and the remaining photon terms the Neumann ones of $X$; the observable depends on the $X$-variables only. The free Wick constant, $\dm_\eps$ and the scheme multiplier are the same on both sides.
\end{proof}

With $Y=X_{J+1}$, $X=\XJ$ this is King's monotonicity in the volume \cite[Thm.~3.2, (3.31)--(3.33)]{K2} between comparison cubes, in the form his proof proves: with one Wick constant for all $J$, the free subtraction King himself switches to at the start of that proof \cite[p.~333]{K2}. With the Dirichlet Wick constant of \cite[(2.13)]{K2}, which increases with $J$, the $c$-number factor moves against the monotone factor and the monotonicity does not follow; that form is not used. The box version with killed sites and the monotonicity of King's RG cubes are in Section~4.

\subsection{Algebra of the scheme, covariance, reflection, continuity}

\begin{lemma}[algebra of the scheme]\label{lem:s02:alg}
Let $a$ be a scheme and $Z^a$ the functional of a state above on a box at spacing $\eps_k$.
\begin{enumerate}[label=\textup{(\alph*)},nosep]
\item For real coefficients $w$ the multiplier $e^{-\sum_jw_ja_k\sS(h_j)}$ is a positive constant depending only on $k$, $w$ and $\sS_{\Lk}(h_j)$. Hence the inequalities of Lemmas~\ref{lem:s02:BFS} and~\ref{lem:s02:cut} between two states on the same lattice with the same sources hold for $Z^a$ in the same direction.
\item[(d)] $|Z^a(\zeta)|\le Z^a(\Re\zeta)$, and $\ln Z^a$ is convex in the real coefficients.
\item[(e)] $Z^a(f,h)\,Z^a(-f,-h)\ge1$ for real $(f,h)$.
\item[(f)] For real $f$ and real $k_1,k_2\ge0$: $Z^a(f,k_1-k_2)\le Z^a(2f,-2k_2)^{1/2}\,Z^a(0,2k_1)^{1/2}$.
\end{enumerate}
\end{lemma}

\begin{proof}
(a) is immediate. (d) $|\Theta(\zeta)|=\Theta(\Re\zeta)$ and the multiplier has modulus $e^{-\sum_j(\Re w_j)a_k\sS(h_j)}$; convexity is H\"older's inequality. (e) $\langle e^{X}\rangle\langle e^{-X}\rangle\ge1$ by Cauchy--Schwarz, and the multipliers cancel. (f) Cauchy--Schwarz gives $\langle e^{F+\Wick{k_1}-\Wick{k_2}}\rangle\le\langle e^{2F-2\Wick{k_2}}\rangle^{1/2}\langle e^{2\Wick{k_1}}\rangle^{1/2}$, and $e^{-a_k\sS(k_1-k_2)}=(e^{2a_k\sS(k_2)})^{1/2}(e^{-2a_k\sS(k_1)})^{1/2}$.
\end{proof}

\begin{lemma}[passive covariance]\label{lem:s02:pass}
Let $O\in\Oth$, $\Lambda$ a lattice of spacing $\eps$ with frame $(f_\alpha)$, and $X$ a box or a torus for $\Lambda$. Equip $O\Lambda$ with the frame $(Of_\alpha)$ and $OX$ with the transported boundary conditions. Then $\langle\Theta(z\,O\cdot f,w\,O\cdot h)\rangle_{O\Lambda,OX}=\langle\Theta(zf,wh)\rangle_{\Lambda,X}$ and $\sS_{O\Lambda}(O\cdot h)=\sS_\Lambda(h)$, likewise for families and for $Z^a$. On a torus the periodic model does not change if the frame $(f_\alpha)$ is replaced by $(\pm f_{\pi(\alpha)})$, $\pi$ a permutation.
\end{lemma}

\begin{proof}
The map $\phi'(Ox)=\phi(x)$, $A'_\alpha(Ox)=A_\alpha(x)$, components taken in the frames $(f_\alpha)$ and $(Of_\alpha)$, is a bijection of the configuration spaces preserving Lebesgue measure; the action is written in the frame, and $\dm_\eps$, $\sW_\eps$, $a_k$ depend on $\eps$ only; plaquettes, centres and field strengths are transported, and $(O\cdot f)_{Oc}(Of_\alpha,Of_\beta)=f_c(f_\alpha,f_\beta)$. A permutation of the frame relabels the components; reversing $f_\alpha$ replaces $A_\alpha(x)$ by $-A_\alpha(x-\eps f_\alpha)$, a bijection of the torus variables preserving the action, which changes the sign of both $F(p)$ and $f_{c(p)}(f_\alpha,f_\beta)$ and preserves the set of plaquette centres.
\end{proof}

\begin{lemma}[the reflection sandwich at fixed spacing, comparison convention]\label{lem:s02:refl}
Let $\sigma_3=\diag(1,1,-1)$ and let $(f,g)\in\cK$ be supported at $\ell^\infty$-distance $>1$ from $\R^3\setminus C_{J-1}$. Then for every $k\ge0$
\[
Z^a_{k,J-1}(\mathrm if,g)\le Z^a_{k,J}(\mathrm i\,\sigma_3\!\cdot\!f,\sigma_3\!\cdot\!g)\le Z^a_{k,J+1}(\mathrm if,g),
\]
and the reversed inequalities hold for $Z^a(f,-g)$.
\end{lemma}

\begin{proof}
$\sigma_3$ preserves $\Lk$ and $\sS_{\Lk}$ and maps $\XJ$ onto $X_+\defeq(\CJ+\eps_ke_3)\cap\Lk$. By Lemma~\ref{lem:s02:pass}, $Z^a_{k,J}(\mathrm i\sigma_3\cdot f,\sigma_3\cdot g)$ is the functional with sources $(f,g)$ of the Neumann--Dirichlet state on $X_+$ in the frame $(e_1,e_2,-e_3)$. Written in the frame $(e_\alpha)$, with $A'_3(x)=-A_3(x-\eps e_3)$, this is the BFS measure $R$ in which $\phi,A_1,A_2$ live on $X_+$ and $A_3$ on $X_+-\eps e_3=\XJ$, with Dirichlet conditions for $\phi$ on $X_+$, photon gradient and mass terms within these sets, and the hopping term in direction $3$ the standard one of the bond $\langle y-\eps e_3,y\rangle$; the plaquette terms $(x;\alpha,3)$ pick up two sign changes and keep their centres, so the observable is $\Theta(\mathrm if,g)$ in the standard frame. $R$ is obtained from the Neumann--Dirichlet state on $X_{J+1}$ by deleting photon gradient terms and hopping cosines and factorising, since $X_+\cup\XJ\subset X_{J+1}$, and the state on $X_{J-1}$ is obtained from $R$ in the same way, since $X_{J-1}\subset\XJ\cap X_+$; the diagonal parts of the cut bonds are the Dirichlet terms of the smaller state, and the sources stay inside $X_{J-1}$. Lemma~\ref{lem:s02:BFS}(a),(b),(e), respectively (d),(e), give the directions, and Lemma~\ref{lem:s02:alg}(a) carries them to $Z^a$.
\end{proof}

\begin{lemma}[continuity of convex functions]\label{lem:s02:baire}
Let $U$ be a non-empty open convex subset of a Fr\'echet space and $\Phi:U\to\R$ a convex Borel function. Then $\Phi$ is continuous on $U$.
\end{lemma}

\begin{proof}
$U$ is a Baire space and $U=\bigcup_{j\in\N}\{\Phi\le j\}$, so some $S\defeq\{\Phi\le j\}$ is non-meagre in $U$; being Borel, $S$ has the Baire property \cite{Kec95}: $O\setminus\mathcal M\subset S$ for a non-empty open $O\subset U$ and a meagre $\mathcal M$. Pick $x_0\in O$ and a balanced open neighbourhood $W$ of $0$ with $x_0+W+W\subset O$. For $w\in W$ and $y\defeq x_0+w$, $y+W\subset O$, and the point reflection $\iota(x)\defeq2y-x$ maps $y+W$ onto itself and preserves meagre sets; so $S\cap(y+W)$ and $\iota(S\cap(y+W))$ are comeagre in $y+W$ and intersect: there is $u\in S$ with $2y-u\in S$, and $\Phi(y)\le\frac12(\Phi(u)+\Phi(2y-u))\le j$. Hence $\Phi\le j$ on $x_0+W$. A convex function on an open convex set bounded above near one point is bounded above near every point $y\in U$ (choose $z\in U$ with $y$ strictly between $x_0$ and $z$, and take the convex combinations of $x_0+W$ with $z$), and a convex function locally bounded above is continuous \cite{RS1}.
\end{proof}
